% GENERATED FILE. Edit canonical lessons, then run npm run generate:books.
% canonical-source-hash: fnv1a64:826fe8eac85b5cc6
% canonical-lessons: LA-C543-cantus, LA-C543-cantio, LA-C543-cantilena, LA-C543-concino, LA-C543-canorus

\chapter{Song}
\label{ch:la-song}
\hlchaptermodality{543}

\begin{chapteropening}
\textbf{By the end of this chapter:} \emph{I can talk about singing and songs.}

Complete the last lesson of chapter 543: I can talk about singing and songs.
\end{chapteropening}

\section[cantus, cantūs]{cantus, cantūs — singing, song; music}

\label{lesson:LA-C543-cantus}



\begin{quote}
Before the new one: say the Latin for to stay in the country, then the Latin for lovely.
\end{quote}

\subsection*{You'll want to know: cantus, cantūs}
\textbf{cantus, cantūs} — \textquotedblleft{}singing, song; music\textquotedblright{}.

\begin{quote}
Cantum avium māne audīmus.
\end{quote}

\begin{quote}
\emph{In the morning we hear the song of the birds.}
\end{quote}

Masculine, fourth declension, from \emph{canō}. Singing, and playing too. English \emph{chant} comes from the verb \emph{cantō}.

\begin{cousinweb}
\textbf{cantus} is the action noun of \emph{canō} \textquotedblleft{}to sing\textquotedblright{}, inherited from a Proto-Indo-European root meaning \textquotedblleft{}to sing\textquotedblright{} (EtymDB). English \emph{chant} belongs to the same family.
\end{cousinweb}

\begin{grammarlens}[title={What you've built}]
One more word for telling, explaining and describing at B1.
\end{grammarlens}

\subsection*{Your turn}
\begin{itemize}
\raggedright
  \item \emph{Say it:} \emph{cantus, cantūs}
  \item \emph{Say it:} \emph{cantus, cantūs}, once more
  \item \emph{Say it:} a sentence of your own with \emph{cantus, cantūs}
  \item \emph{Recall:} say the Latin for to stay in the country, then the Latin for lovely, then say \emph{cantus, cantūs} again
\end{itemize}

\begin{tcolorbox}[breakable,colback=teal!4,colframe=teal!35!black,title={Before you move on}]
What does cantus, cantūs mean? (\textquotedblleft{}Singing, song; music\textquotedblright{}.) Say it once more.
\end{tcolorbox}
\section[cantiō, cantiōnis]{cantiō, cantiōnis — a song; a spell}

\label{lesson:LA-C543-cantio}



\begin{quote}
Before the new one: say the Latin for lovely, then the Latin for singing.
\end{quote}

\subsection*{You'll want to know: cantiō, cantiōnis}
\textbf{cantiō, cantiōnis} — \textquotedblleft{}a song; a spell\textquotedblright{}.

\begin{quote}
Ancillae in culīnā cantiōnem canunt.
\end{quote}

\begin{quote}
\emph{The maids sing a song in the kitchen.}
\end{quote}

Feminine, third declension, from \emph{canō}. Also a magic spell sung over someone. French \emph{chanson} and Italian \emph{canzone} come from it.

\begin{cousinweb}
\textbf{cantiō} is the action noun of \emph{canō} \textquotedblleft{}to sing\textquotedblright{} (Lewis \& Short; EtymDB). French \emph{chanson} and Italian \emph{canzone} come from it.
\end{cousinweb}

\begin{grammarlens}[title={What you've built}]
One more word for telling, explaining and describing at B1.
\end{grammarlens}

\subsection*{Your turn}
\begin{itemize}
\raggedright
  \item \emph{Say it:} \emph{cantiō, cantiōnis}
  \item \emph{Say it:} \emph{cantiō, cantiōnis}, once more
  \item \emph{Say it:} a sentence of your own with \emph{cantiō, cantiōnis}
  \item \emph{Recall:} say the Latin for lovely, then the Latin for singing, then say \emph{cantiō, cantiōnis} again
\end{itemize}

\begin{tcolorbox}[breakable,colback=teal!4,colframe=teal!35!black,title={Before you move on}]
What does cantiō, cantiōnis mean? (\textquotedblleft{}A song; a spell\textquotedblright{}.) Say it once more.
\end{tcolorbox}
\section[cantilēna, cantilēnae]{cantilēna, cantilēnae — an old song, the same old tune}

\label{lesson:LA-C543-cantilena}



\begin{quote}
Before the new one: say the Latin for singing, then the Latin for a song.
\end{quote}

\subsection*{You'll want to know: cantilēna, cantilēnae}
\textbf{cantilēna, cantilēnae} — \textquotedblleft{}an old song, the same old tune\textquotedblright{}.

\begin{quote}
Avia eandem cantilēnam semper canit.
\end{quote}

\begin{quote}
\emph{Grandmother always sings the same old song.}
\end{quote}

Feminine, first declension. \emph{Cantilēnam eandem canis} — you sing the same old tune: you keep repeating yourself.

\begin{cousinweb}
\textbf{cantilēna} is built on \emph{cantillō} \textquotedblleft{}to hum\textquotedblright{}, a diminutive of \emph{cantō} \textquotedblleft{}to sing\textquotedblright{} (Lewis \& Short).
\end{cousinweb}

\begin{grammarlens}[title={What you've built}]
One more word for telling, explaining and describing at B1.
\end{grammarlens}

\subsection*{Your turn}
\begin{itemize}
\raggedright
  \item \emph{Say it:} \emph{cantilēna, cantilēnae}
  \item \emph{Say it:} \emph{cantilēna, cantilēnae}, once more
  \item \emph{Say it:} a sentence of your own with \emph{cantilēna, cantilēnae}
  \item \emph{Recall:} say the Latin for singing, then the Latin for a song, then say \emph{cantilēna, cantilēnae} again
\end{itemize}

\begin{tcolorbox}[breakable,colback=teal!4,colframe=teal!35!black,title={Before you move on}]
What does cantilēna, cantilēnae mean? (\textquotedblleft{}An old song, the same old tune\textquotedblright{}.) Say it once more.
\end{tcolorbox}
\section[concinō, concinere]{concinō, concinere — to sing together, sound in harmony}

\label{lesson:LA-C543-concino}



\begin{quote}
Before the new one: say the Latin for a song, then the Latin for an old song.
\end{quote}

\subsection*{You'll want to know: concinō, concinere}
\textbf{concinō, concinere} — \textquotedblleft{}to sing together, sound in harmony\textquotedblright{}.

\begin{quote}
Puellae in templō concinunt.
\end{quote}

\begin{quote}
\emph{The girls sing together in the temple.}
\end{quote}

\emph{concinō, concinere, concinuī}: \emph{con-} + \emph{canō}.

\begin{cousinweb}
\textbf{concinō} is built on \emph{con-} \textquotedblleft{}together\textquotedblright{} and \emph{canō} \textquotedblleft{}to sing\textquotedblright{} (Lewis \& Short; EtymDB).
\end{cousinweb}

\begin{grammarlens}[title={What you've built}]
One more word for telling, explaining and describing at B1.
\end{grammarlens}

\subsection*{Your turn}
\begin{itemize}
\raggedright
  \item \emph{Say it:} \emph{concinō, concinere}
  \item \emph{Say it:} \emph{concinō, concinere}, once more
  \item \emph{Say it:} a sentence of your own with \emph{concinō, concinere}
  \item \emph{Recall:} say the Latin for a song, then the Latin for an old song, then say \emph{concinō, concinere} again
\end{itemize}

\begin{tcolorbox}[breakable,colback=teal!4,colframe=teal!35!black,title={Before you move on}]
What does concinō, concinere mean? (\textquotedblleft{}To sing together, sound in harmony\textquotedblright{}.) Say it once more.
\end{tcolorbox}
\section[canōrus, canōra, canōrum]{canōrus, canōra, canōrum — melodious, tuneful}

\label{lesson:LA-C543-canorus}



\begin{quote}
Before the new one: say the Latin for an old song, then the Latin for to sing together.
\end{quote}

\subsection*{You'll want to know: canōrus, canōra, canōrum}
\textbf{canōrus, canōra, canōrum} — \textquotedblleft{}melodious, tuneful\textquotedblright{}.

\begin{quote}
Vōx puellae canōra est.
\end{quote}

\begin{quote}
\emph{The girl's voice is melodious.}
\end{quote}

From \emph{canor} — melody, and so from \emph{canō}.

\begin{cousinweb}
\textbf{canōrus} is built on \emph{canor} \textquotedblleft{}melody\textquotedblright{}, from \emph{canō} \textquotedblleft{}to sing\textquotedblright{} (Lewis \& Short; EtymDB).
\end{cousinweb}

\begin{grammarlens}[title={What you've built}]
One more word for telling, explaining and describing at B1.
\end{grammarlens}

\subsection*{Your turn}
\begin{itemize}
\raggedright
  \item \emph{Say it:} \emph{canōrus, canōra, canōrum}
  \item \emph{Say it:} \emph{canōrus, canōra, canōrum}, once more
  \item \emph{Say it:} a sentence of your own with \emph{canōrus, canōra, canōrum}
  \item \emph{Recall:} say the Latin for an old song, then the Latin for to sing together, then say \emph{canōrus, canōra, canōrum} again
\end{itemize}

\begin{tcolorbox}[breakable,colback=teal!4,colframe=teal!35!black,title={Before you move on}]
What does canōrus, canōra, canōrum mean? (\textquotedblleft{}Melodious, tuneful\textquotedblright{}.) Say it once more.
\end{tcolorbox}
